\section*{Chapitre \ref{ch:rc:vanilla-rates-options} --- Options de taux vanille}

\begin{solution}{exo:rc:vanilla-rates-options:1}
$10\,000\,000\times(0.034-0.030)\times0.5069 = \text{EUR}~20\,276$, versé à la
fin de la période, six mois après la fixation.
\end{solution}

\begin{solution}{exo:rc:vanilla-rates-options:2}
$\sigma_B\approx\sigma_N/F = 80/250 = 32~\%$~; la valeur exacte 32.7~\% est légèrement
plus élevée car $2\Phi(x/2)-1 < x/\sqrt{2\pi}$ pour $x = \sigma_B\sqrt T > 0$.
\end{solution}

\begin{solution}{exo:rc:vanilla-rates-options:3}
Le taux de la première période est déjà fixé à la date de transaction~: son caplet n'a
plus aucune optionnalité, seulement un paiement connu~; il est donc écarté et l'optionnalité
du cap commence avec la deuxième période.
\end{solution}

\begin{solution}{exo:rc:vanilla-rates-options:4}
Ci-dessus~: la valeur du cap de deux ans à 66 points de base inclut les caplets de la
première année, valant 58~; pour que la moyenne atteigne 66, les caplets de la deuxième
année doivent être plus élevés. Le stripping donne 66.6.
\end{solution}

\begin{solution}{exo:rc:vanilla-rates-options:5}
Le prix physique multiplie l'option par l'annuité collatéralisée 7.740~;
la formule du taux au pair utilise $P(0,T)a(S_0) = 7.579$, qui actualise la
annuité au taux de swap Euribor 3.098~\%, supérieur aux taux d'actualisation €STR~:
2.1~\% de moins. Le prix comptant collatéralisé se règle à la valeur du swap sur la
\omterm{def:rc:multi-curve-and-collateral-discounting:curves}{courbe d'actualisation} de la chambre de compensation, ce que livre la swaption physique,
donc les deux prix coïncident.
\end{solution}

\begin{solution}{exo:rc:vanilla-rates-options:6}
Dans le modèle normal, un mouvement de tant de points de base est aussi probable quel que
soit le niveau~; par rapport à un strike bas, c'est un grand mouvement en pourcentage, par rapport à un
strike élevé, un petit, si bien que la \omterm{def:rc:vanilla-rates-options:lognormal}{volatilité lognormale} implicite baisse avec le
strike. Un marché lognormal plat ressemble à un skew normal positif~: la volatilité normale
$\approx\sigma_B\sqrt{FK}$ croît avec le strike.
\end{solution}

\begin{solution}{exo:rc:vanilla-rates-options:7}
1.97~\%~: la vente du floor à cinq ans à 1.97~\% finance le cap à 3~\%. La
société paie alors au moins 1.97~\% plus sa marge, même si l'Euribor baisse
encore. (Utiliser la \omterm{def:rc:vanilla-rates-options:flat}{volatilité de caplet} à 3~\% pour tous les strikes revient à ignorer le smile.)
\end{solution}

\begin{solution}{exo:rc:vanilla-rates-options:8}
76 points de base est la \omterm{def:rc:vanilla-rates-options:flat}{volatilité plate} de l'ensemble du cap à cinq ans, une moyenne
sur tous ses caplets~; les caplets dont la fixation intervient après quatre ans ont leur
propre volatilité issue du stripping, 78.0. À 76, elles valent 886\,956~EUR sur un notionnel de 200~millions d'euros,
à 78.0, 916\,329~EUR~: la sous-évaluation s'élève à 29\,373~EUR.
\end{solution}

\begin{solution}{pb:rc:vanilla-rates-options:1}
\textbf{1.} Neuf~: dix périodes de six mois, moins la première, déjà fixée.
\textbf{2.} 2\,157\,482~EUR, soit 1.079~\% du notionnel.
\textbf{3.} Divisé par l'annuité des paiements de caplets (4.286)~: 25.2 points de base
par an.
\textbf{4.} Le même montant, 2\,157\,482~EUR~: les volatilités issues du stripping ont été calibrées de
façon que le cap à cinq ans se réévalue à sa \omterm{def:rc:vanilla-rates-options:flat}{volatilité plate}.
\textbf{5.} Les forwards sont des espérances~; avec une volatilité normale de 66 à 79 points
de base par an, des fixations supérieures à 3~\% dans les années suivantes sont tout à fait
probables, et le cap paie dans ces états.
\textbf{6.} 43\,093~EUR par point de base.
\textbf{7.} Environ $10\times43\,093 = \text{EUR}~430\,929$.
\textbf{8.} Croît~: un point de \omterm{def:rc:vanilla-rates-options:lognormal}{volatilité lognormale} vaut environ $F$ points de
volatilité normale ($\sigma_N\approx F\sigma_B$), donc la vega lognormale vaut environ
$F$ fois la vega normale, qui ne dépend pas du niveau~; des taux plus élevés,
une vega lognormale plus grande. La vega normale est le chiffre le plus stable à publier.
\textbf{9.} La banque qui l'a vendu~; elle se couvre en achetant des caps ou des caplets sur le
marché, ou des swaptions pour les échéances plus longues, et gère le risque résiduel.
\textbf{10.} 1\,827\,377~EUR~: à 25~\% plate, les caplets sur des forwards proches de
2.4~\% ont des volatilités normales d'environ 60 points de base, inférieures à celles du marché~;
la cotation lognormale sous-évaluerait le cap de 15~\%.
\textbf{11.} $200\text{ millions d'euros}\times0.005\times\sum\delta_k$ avec $\sum\delta_k =
4.569$~: 4\,569\,444~EUR.
\textbf{12.} 4\,285\,872~EUR ($0.005\times200\text{ millions d'euros}\times$ annuité 4.286).
\textbf{13.} Environ 3.25~\%~: la prime représente 25.2 points de base par an du
notionnel, donc les fixings doivent en moyenne dépasser le strike de 25 points de base.
\textbf{14.} Assurance~: protection contre les états où les taux ont monté,
qui ne se sont pas réalisés~; la prime est le prix de cette protection.
\textbf{15.} Chaque caplet paie sur sa propre fixation, si bien qu'une fixation élevée dans une période
n'est pas compensée par une fixation basse dans une autre~; une option sur la moyenne serait moins chère.
\textbf{16.} Le swap fige le taux au forward (environ 2.4~\% à 3~\% sur les
périodes) sans coût initial et supprime le bénéfice de la baisse des taux~; le cap
coûte 25 points de base par an et le conserve.
\textbf{17.} 1.97~\%~: la société paie son cap en renonçant au bénéfice des
fixations inférieures à 1.97~\%.
\textbf{18.} Écart cours acheteur--cours vendeur sur la volatilité qu'elle doit acheter pour se couvrir, la charge
de crédit d'une entreprise qui ne poste pas de collatéral (chapitres~18 et~19), et le
capital que consomme la transaction.
\textbf{19.} Résultat nommé~: \emph{le cap du trésorier} coûte 2\,157\,482~EUR (1.079~\%
du notionnel, 25.2 points de base par an) avec une vega de 43\,093~EUR par point de base
de volatilité normale.
\textbf{20.} Protection contre les taux élevés tout en conservant le bénéfice des taux bas.
\end{solution}

\begin{solution}{iq:rc:vanilla-rates-options:1}
Une option versant $N\delta\max(F-K,0)$ à la fin d'une période sur le taux fixé à
son début. Son forward $F_k(t) = (P(t,T_{k-1})/P(t,T_k)-1)/\delta$ est un actif
négocié divisé par $P(t,T_k)$, c'est donc une martingale sous la probabilité forward
$T_k$, ce qui fait que le prix du caplet vaut $N\delta P(0,T_k)$ fois un call sur une
martingale.
\iqlookfor{l'argument du numéraire.}
\end{solution}

\begin{solution}{iq:rc:vanilla-rates-options:2}
À la monnaie, Bachelier donne $\sigma_N\sqrt{T}\varphi(0) = \sigma_N\sqrt{T/2\pi}$
et Black donne $F(2\Phi(\sigma_B\sqrt T/2)-1)\approx F\sigma_B\sqrt{T}\varphi(0)$.
En égalisant~: $\sigma_N\approx F\sigma_B$~; plus généralement $\sigma_N\approx\sigma_B\sqrt{FK}$
pour des strikes voisins.
\iqlookfor{le développement et la règle en $\sqrt{FK}$.}
\end{solution}

\begin{solution}{iq:rc:vanilla-rates-options:3}
On price chaque cap cité à sa \omterm{def:rc:vanilla-rates-options:flat}{volatilité plate}~; avec des caps d'échéance croissante,
on résout pour une volatilité des nouveaux caplets telle que le cap plus long se reprice,
les caplets antérieurs étant conservés (ou on ajuste une courbe paramétrique lisse à l'ensemble des caps en une fois).
Problèmes~: volatilités négatives ou oscillantes lorsque les cotations sont incohérentes ou
que les échéances sont proches, sensibilité du strip à une seule mauvaise cotation, un strip
différent pour chaque strike (le smile), et la convention du premier caplet.
\iqlookfor{la résolution séquentielle et ses modes de défaillance.}
\end{solution}

\begin{solution}{iq:rc:vanilla-rates-options:4}
La volatilité normale est plus stable lorsque les taux bougent, puisque les variations de taux
en points de base dépendent peu du niveau~; elle fonctionne avec des taux négatifs~; les vegas
normaux s'additionnent entre instruments de forwards différents~; et c'est elle que le marché cote, si bien que les
ratios de couverture correspondent aux cotations.
\iqlookfor{stabilité du smile et additivité.}
\end{solution}

\begin{solution}{iq:rc:vanilla-rates-options:5}
Le payoff au rendement au pair $a(S_T)(S_T-K)^+$ n'est pas le numéraire d'annuité multiplié par le
payoff, si bien que la valorisation sous la probabilité annuité ne s'applique pas~; la formule
$P(0,T)a(S_0)\,\mathrm{Black}$ ignore la convexité de $a(S)$ et mélange
des taux d'actualisation, et elle n'est pas sans arbitrage par rapport aux swaptions physiques. Le
marché (swaptions en euros, novembre 2018) a fait passer la valeur par défaut au
prix cash collatéralisé, qui se règle à la valeur du swap sur la courbe de la chambre de compensation.
\iqlookfor{payoff face au numéraire, et la correction du marché.}
\end{solution}

\begin{solution}{iq:rc:vanilla-rates-options:6}
La formule de Black prend $\ln(F/K)$~: un forward ou un strike négatif produit un NaN ou
une exception. On valorise dans le modèle normal ou dans un modèle lognormal décalé (chapitre
5), on convertit les cotations de volatilité en conséquence, et on teste le pricer sur des
forwards et des strikes négatifs.
\iqlookfor{le logarithme et le changement de modèle.}
\end{solution}
